\documentclass[11pt,a4paper]{article}
\usepackage[T1]{fontenc}
\usepackage[utf8]{inputenc}
\usepackage{amsmath,amssymb,xcolor}
\usepackage[margin=25mm]{geometry}
\usepackage[colorlinks=true,urlcolor=blue]{hyperref}
\definecolor{nexorablue}{HTML}{2459B5}
\setlength{\parindent}{0pt}
\setlength{\parskip}{7pt}
\setlength{\emergencystretch}{2em}
\renewcommand{\familydefault}{\sfdefault}
\begin{document}

{\small\bfseries\color{nexorablue} NEXORA CPGE}\par
{\small MP / NEX-TOP-01 / Corrigé}\par
{\LARGE\bfseries Un losange incliné}\par
\vspace{0.25em}\hrule\vspace{0.6em}

\section*{Norme et constantes}
La transformation $(x,y)\mapsto(u,v)$ est linéaire et inversible. Elle transporte la norme $1$. On a $N(x,y)\le |x|+3|y|\le4\|(x,y)\|_\infty$. De même $|y|\le N$ et $|x|=|u-2v|\le2(|u|+|v|)=2N$.

\section*{Optimalité et sommets}
Les quotients $N/\|\cdot\|_\infty$ valent $4$ en $(1,1)$ et $1/2$ en $(-2,1)$. Les meilleures constantes sont donc $a=1/2$ et $b=4$. Les sommets du losange $|u|+|v|\le1$ deviennent $(1,0)$, $(-1,0)$, $(-2,1)$ et $(2,-1)$.


\end{document}
